\scanpage{65}
\deleted{\unsure{équivaut} à celle \unsure{de} \unsure{nous} \unsure{on|en|un} \unsure{peut|faut|fct} \unsure{revenir}}
\[
\boxed{\underline{\lambda}^{g}\lambda^{-1}=}
\]

Il y a une solution triviale $\lambda=\frac{1}{L}$ i.e. $L'(t)=1$. Mais nous \unsure{désirons|devons} \deleted{\unsure{que}} \inserted{imposer} que $\lambda$ \deleted{\unsure{$L$}} soit de ``niveau'' \unsure{strictmt|strict.} inf à $L$.

Pour préciser ce point, nous \unsure{supposons},

\deleted{$L(0)=1$, que les racines et pôles \unsure{de} $L$ sont en valeurs absolues de la forme $p^{-\frac{i}{2}}$ avec $\deleted{2}\rho\leq i\leq\deleted{2}\rho+2n$ i.e. $\rho\leq\frac{i}{2}\leq\rho+n$, \deleted{et \unsure{modulo multiplication|multipliant|multipliées|modulo coefficients} par une \unsure{puissance} de $p^{1/2}$, \unsure{des|dans} des valeurs \unsure{propres|propre} \unsure{entières}} et on veut se restreindre aux $\lambda$ ayant \uline{la $=$ propriété}, \deleted{et de plus} (de poids compris entre $\rho$ et $\rho+n$) et de plus de ``niveau'' $\leq\deleted{\unsure{\rho+n+1}}$}
% margin note written sideways
\margin{et sont \unsure{tous|toutes} de la forme $(p^{1/2})^{\deleted{\unsure{i}}\unsure{\nu}}\zeta_i$, où $\nu$ un entier $\zeta_i$ un \uline{entier alg.}}
\[
L(t)=\underbrace{L_{\rho}(t)\cdots L_{\rho\deleted{\unsure{b}}+n}(t)}_{\unsure{\longrightarrow|1|\ne}}L'_{\rho+n-1}(pt)L'_{\rho+n-2}(p^2t)\cdots L'_{\rho}(p^nt)
\]
où les valeurs propres inverses des $L_{\deleted{\unsure{2}}i}$ sont \deleted{\unsure{et des}} des entiers alg en val. absolue $p^{i/2}$, et de $=$ pour $L'_i$ (donc $L'_{\rho+n-i}(p^it)$ \unsure{à|a} des valeurs propres inverses qui sont des entiers alg de val. absolue $p^{\frac{\rho+n-i}{2}}p^i=p^{\frac{\rho+n+i}{2}}$.
% archive sheet number: 64
